\begingroup
\ArithmeticHead{245}{R}{GAUSS AN SCHUMACHER.}
\SourceLeaf{252b}
\pdfbookmark[1]{[6.] Schumacher an Gauss, 5. April 1816}{x1-correspondence-6}
\NumberedItem{6}
\begin{center}\textsc{Schumacher} an \textsc{Gauss}. Kopenhagen, d. $5^{\mathrm{ten}}$ April 1816.\end{center}
\HistoricCitation{[Briefwechsel zwischen \textsc{Gauss} und \textsc{Schumacher} I, Altona 1860, S. 123—124.]}
\noindent .\enspace .\enspace .\enspace .\quad \{Ein Landsmann von Ihnen ist hier Professor der Mathematik. Er heisst \textsc{Degen}. Vor einiger Zeit sagte er mir, er beschäftige sich schon seit-
\clearpage
\ArithmeticHead{246}{L}{ANALYSIS. BRIEFWECHSEL.}
\SourceLeaf{253}
\noindent dem er \textsc{Tetens} Schüler sei, mit gewissen Mediis arithmetico-geometricis, die von den Ihrigen verschieden sind. Er nimmt nemlich zwei Grössen, von denen $a$ die grössere, $b$ die kleinere seyn soll (die Stelle ist gleichgültig)
\[
\begin{array}{llll}
a,&\tfrac12(a+b)=a',&\tfrac12(a'+b')=a'',&\text{u. s. w.}\\[4pt]
b,&\surd ab=b',&\surd a'b'=b'',&\text{u. s. w.}
\end{array}
\]
\par\indent Er selbst weiss nur, dass diese Reihe einen Limes hat, den man bald findet. Ich habe mich sehr mit diesem Limes beschäftigt und bin jetzt so glücklich gewesen, einen Zusammenhang mit der Ellipse zu finden, den ich mir die Freiheit nehme Ihnen vorzulegen. Da der Limes dieser Reihe so wie ich ihn dargestellt habe, von $\dfrac{a}{b}$ abhängt, so bezeichne ich ihn mit $M\left(\dfrac{b}{a}\right)$. Ich nenne ferner die sehr schnell convergirende Reihe (aus den Differenzen der successiven $a,a',a''\ldots$ gebildet)
\[
\frac{a-a'}{2a'}+\frac{a-a'}{2a'}\, .\,\frac{a'-a''}{2a''}
+\frac{a-a'}{2a'}\, .\,\frac{a'-a''}{2a''}\, .\,\frac{a''-a'''}{2a'''}+\ldots
\]
\noindent $\mathfrak{M}\left(\dfrac{b}{a}\right)$, dann ist
\[
\frac12\pi\left(\frac{a'}{M\left(\dfrac{b}{a}\right)}-\frac{(a-a')\mathfrak{M}\left(\dfrac{b}{a}\right)}{M\left(\dfrac{b}{a}\right)}\right)
=\text{Quadrans Ellipseos}\left(\frac{b'}{a'}\right),
\]
\noindent wo das Zeichen $\left(\dfrac{b'}{a'}\right)$ eine Ellipse bedeutet, deren grosse Axe $=1$, kleine Axe $=\dfrac{b'}{a'}$ ist. Ich habe nun \textsc{Degen} angezeigt, ich habe seinen Limes gefunden und halte es für billig zu warten, ob er dasselbe finde. Wie ich vor 14 Tagen an \textsc{Lindenau} schrieb, hatte ich den Limes noch nicht.

\par\indent Gewiss haben Sie bei Ihren Mediis auch dieses angesehen, und Sie würden mich sehr durch ein paar Worte darüber verbinden.

\par\indent Leben Sie wohl, werthester Freund und vergessen Sie nicht
\begin{flushright}Ihnen ganz eignen\par\medskip\textsc{Schumacher.}\}\end{flushright}
\clearpage
\endgroup
